\section{Introduction}

Finite memory under drift has two failure modes. One is distributional: a
window that is too long becomes temporally stale, so the retained evidence no
longer tracks the present distribution accurately. The other is sequential: a
detector run on the same stream may still remain below alarm threshold while the
window is already invalid. These two failures need not occur on the same
timescale.

The temporal-validity horizon identifies the first clock. In finite-memory
distribution tracking under drift, the risk balance
\[
  \mathbb E\,d(\widehat P_t^{(n)},P_t) \le C_K n^{-a} + C_S \zeta n^H
\]
induces a horizon scale $n_t^* \asymp (C_K/\zeta_t)^{1/(a+H)}$. For a fixed
operating memory level $n$, validity expires once the local horizon contracts
below $n$. That defines the expiry time $\tau_{\mathrm{valid}}(n)$.

Sequential detection introduces a second clock. A stopping time
$\tau_{\mathrm{detect}}$ is designed to react to change while satisfying a
false-alarm constraint. The central object of this paper is therefore the
validity-detection delay
\[
  \Delta_{\mathrm{val-det}} \coloneqq \tau_{\mathrm{detect}} - \tau_{\mathrm{valid}}(n).
\]
Positive values mean that retained evidence has already become invalid before a
detector alarm appears.

That positive delay is structural rather than a tuning accident. The detector-
class sign law, and the lower-bound problem that follows it, show that under
false-alarm control there exist H\"older drift families for which detection must
lag validity expiry with nontrivial probability.

The benchmark proof model is one-dimensional Gaussian location with deterministic
H\"older mean path and fixed observation noise. That model separates the
finite-memory horizon inherited from distribution tracking, the local post-
expiry evidence available to the detector, and the false-alarm constraint that
limits how aggressively stopping rules can react. Detector families enter as
benchmark instantiations of a broader stopping-time class.

The empirical signal already points in that direction. On a compact
false-alarm-calibrated frontier over roughness exponents
$H\in\{0.5,0.75,1.0\}$, observation-based ADWIN, Page--Hinkley, and CUSUM keep
positive validity-detection gaps across the tested grid, while residual-input
variants and KSWIN often trade larger delays for missed alarms. Those results do
not prove the delay law, but they identify the right theorem target.
